%----------------------------------------------------------------------------------------
%	PACKAGES AND OTHER DOCUMENT CONFIGURATIONS
%----------------------------------------------------------------------------------------

\documentclass[10pt,stdletter,dateno,sigleft]{newlfm} % Extra options: 'sigleft' for a left-aligned signature, 'stdletternofrom' to remove the from address, 'letterpaper' for US letter paper - consult the newlfm class manual for more options

%\usepackage{charter} % Use the Charter font for the document text
\usepackage{graphics}

\usepackage{standalone}

%% with this patch, the newlfm should work (BEGIN)
\usepackage{etoolbox}
\makeatletter
\patchcmd{\@zfancyhead}{\fancy@reset}{\f@nch@reset}{}{}
\patchcmd{\@set@em@up}{\f@ncyolh}{\f@nch@olh}{}{}
\patchcmd{\@set@em@up}{\f@ncyolh}{\f@nch@olh}{}{}
\patchcmd{\@set@em@up}{\f@ncyorh}{\f@nch@orh}{}{}
\makeatother
%% with this patch, the newlfm should work (END)


\newsavebox{\Luiuc}\sbox{\Luiuc}{\parbox[b]{1.75in}{\vspace{0.5in}
\includegraphics[width=1.3\linewidth]{/Users/hectorbahamonde/Job_Market/logo.png}}} % Company/institution logo at the top left of the page
\makeletterhead{Uiuc}{\Lheader{\usebox{\Luiuc}}}

\newlfmP{sigsize=50pt} % Slightly decrease the height of the signature field
\newlfmP{addrfromphone} % Print a phone number under the sender's address
\newlfmP{addrfromemail} % Print an email address under the sender's address
\PhrPhone{\texttt{p}} % Customize the "Telephone" text
\PhrEmail{\texttt{e}} % Customize the "E-mail" text

\lthUiuc % Print the company/institution logo

% email input begin
\newread\fid
\newcommand{\readfile}[1]% #1 = filename
{\bgroup
  \endlinechar=-1
  \openin\fid=#1
  \read\fid to\filetext
  \loop\ifx\empty\filetext\relax% skip over comments
    \read\fid to\filetext
  \repeat
  \closein\fid
  \global\let\filetext=\filetext
\egroup}
\readfile{/Users/hectorbahamonde/Bibliografia_PoliSci/email.txt}
% email input end

\usepackage{color} % for colors

\usepackage{hyperref}
\hypersetup{
    %bookmarks=true,         % show bookmarks bar?
    unicode=false,          % non-Latin characters in Acrobat’s bookmarks
    pdftoolbar=true,        % show Acrobat’s toolbar?
    pdfmenubar=true,        % show Acrobat’s menu?
    pdffitwindow=true,     % window fit to page when opened
    pdfstartview={FitH},    % fits the width of the page to the window
    pdftitle={My title},    % title
    pdfauthor={Author},     % author
    pdfsubject={Subject},   % subject of the document
    pdfcreator={Creator},   % creator of the document
    pdfproducer={Producer}, % producer of the document
    pdfkeywords={keyword1} {key2} {key3}, % list of keywords
    pdfnewwindow=true,      % links in new window
    colorlinks=true,       % false: boxed links; true: colored links
    linkcolor=blue,          % color of internal links (change box color with linkbordercolor)
    citecolor=blue,        % color of links to bibliography
    filecolor=blue,      % color of file links
    urlcolor=blue           % color of external links
}

\hypersetup{
  colorlinks = true,
  urlcolor = blue,
  pdfpagemode = UseNone
}


%%% bib begin
\usepackage[american]{babel}
\usepackage{csquotes}
%\usepackage[style=chicago-authordate,doi=false,isbn=false,url=false,eprint=false]{biblatex}

\usepackage[authordate,isbn=false,doi=false,url=false,eprint=false]{biblatex-chicago}
\DeclareFieldFormat[article]{title}{\mkbibquote{#1}} % make article titles in quotes
\DeclareFieldFormat[thesis]{title}{\mkbibemph{#1}} % make theses italics

\AtEveryBibitem{\clearfield{month}}
\AtEveryCitekey{\clearfield{month}}

\addbibresource{/Users/hectorbahamonde/Bibliografia_PoliSci/library.bib} 
\addbibresource{/Users/hectorbahamonde/Bibliografia_PoliSci/Bahamonde_BibTex2013.bib} 

% USAGES
%% use \textcite to cite normal
%% \parencite to cite in parentheses
%% \footcite to cite in footnote
%% the default can be modified in autocite=FOO, footnote, for ex. 
%%% bib end

%----------------------------------------------------------------------------------------
%	YOUR NAME AND CONTACT INFORMATION
%----------------------------------------------------------------------------------------

%\namefrom{\vspace{-3cm}Hector Bahamonde, PhD} % Name

\addrfrom{
{\vspace{-2cm}}\\ % Date
[12pt]
{\bf Hector Bahamonde, PhD}\\
{\color{blue}{\bf Senior Researcher}}\\
\emph{{\color{blue}INVEST}}\\
{\color{blue}University of Turku}\\
Turku, Finland
}

\phonefrom{+358 (0) 40-362-1438} % Phone number

\emailfrom{\href{mailto:hibano@utu.fi}{hibano@utu.fi} \\ 
\texttt{w}: \href{http://www.hectorbahamonde.com}{www.HectorBahamonde.com}\\
\today}% Link to last version


%----------------------------------------------------------------------------------------
%	ADDRESSEE AND GREETING/CLOSING
%----------------------------------------------------------------------------------------

%\greetto{\vspace{-1cm}Dear Members of the Search Committee,} % Greeting text
%\closeline{\vspace{-1cm}\includegraphics[width=3cm]{/Users/hectorbahamonde/Administracion/signature.pdf} \hspace{-4cm}} % Closing text




%----------------------------------------------------------------------------------------

\begin{document}
\begin{newlfm}

%----------------------------------------------------------------------------------------
%	LETTER CONTENT
%----------------------------------------------------------------------------------------

%\vspace{-1cm}
\vspace{-2cm}I am writing to apply for the tenure-track position available at the Faculty of Social Sciences of the University of Turku. Currently, I am a Senior Researcher at the INVEST flagship unit, University of Turku. In this cover letter I will outline my motivations to the task and vision for the future development of research, teaching and societal interaction in the discipline.

First, I'll briefly describe my trajectory. Before joining University of Turku in 2021, I was an Assistant Professor (tenure-track) for three years at Universidad de O'Higgins (UOH), Chile. Previously, I had served one year as a post-doctoral fellow at Tulane University (New Orleans, Louisiana, U.S.). I received my PhD in Political Science from Rutgers University (New Brunswick, New Jersey, U.S.) in the spring of 2017. {\bf My research and teaching interests are the political economy of inequality and democratic development, as well as quantitative/statistical, experimental and quasi-experimental methods}. 




% bio and my research in general
%I obtained my B.A. in Political Science at Catholic University of Chile. Then I received a full fellowship from Rutgers University-New Brunswick, NJ, United States, to pursue my PhD in Political Science, which I earned in May 2017. I spent one year as a post-doctoral fellow at Tulane University-New Orleans, LA, U.S. {\bf My research explores the political economy of inequality and democratic development. I am particularly passionate and versed in quantitative/statistical, experimental and quasi-experimental methods}. 

% Experience at INVEST (general)
I've been lucky enough to be a Senior Researcher at INVEST since September 2021. This March 2022 will mark 7 months of hard-work which come with a profound sense of intellectual stimulation and professional accomplishment. In this period of time, I've published ``\href{https://link.springer.com/article/10.1057/s41269-020-00174-4}{Still for Sale: The Micro-Dynamics of Vote Selling in the United States, Evidence from a List Experiment}'' (Acta Politica, \emph{A1}) and ``\href{https://www.sciencedirect.com/science/article/pii/S0176268021000471?via%3Dihub}{Inclusive Institutions, Unequal Outcomes: Democracy, State Capacity, and Income Inequality}'' (European Journal of Political Economy, \emph{A1}). But just as important, I've enjoyed so much the process of collaborating with different researchers from INVEST, the Turku School of Economics and the Political Science Unit at the University of Turku. Not only that, I am actively collaborating with scholars from the \AA bo Akademi who share similar research interests with me---I explain these collaboration instances further below. In sum, I have embraced with joy this intellectual community which has always been welcoming, warm and positively challenging to me. As I explain below, the tenure-track position would allow me to continue contributing even more to the University of Turku.

% Strengthen 
This position would allow me to strengthen even more the methodological and substantive capabilities of the Faculty of Social Sciences and the INVEST ecosystem. In particular, the post will allow me to contribute with research and teaching outputs centered around experimental and quasi-experimental designs, as well as, quantitative/statistical methods. This contribution will not only materialize in peer-reviewed papers, but also in the classroom, with service, by obtaining external funds, and by joining the community outside the walls of our campus (e.g. societal interaction). Let me elaborate some more. 

% Research: Experimental and quasi-experimental
Experimental and quasi-experimental social science is very much what I have always done, and continue using in my research and teaching in the classroom. While randomization is the gold standard to establish causal relationships, quasi-experimental designs---if certain assumptions hold---often times offer good-enough second-bests options. For instance, using matching techniques I published in 2018 ``\href{https://journals.sub.uni-hamburg.de/giga/jpla/article/view/1121/1128}{Aiming Right at You: Group versus Individual Clientelistic Targeting in Brazil}'' (Journals of Politics in Latin America, \emph{A1}). In this paper I argued that vote-buying was not only a function of earnings but also about income inequality. However, every time that's possible---and ethically feasible---experiments should be conducted. Following this reasoning, we developed an economic experiment for the ``\href{https://github.com/hbahamonde/Economic_Experiment_Vote_Selling/raw/master/Vote_Selling_Bahamonde_Canales_Paper.pdf}{Electoral Risk and Vote Buying, Introducing Prospect Theory in the Experimental Study of Clientelism}'' paper. In this manuscript we find that vote buying is most likely when parties are probable \emph{winners} (not \emph{losers}, as traditional research explains). The paper just received  {\bf a revise and resubmit} in {\bf \emph{Electoral Studies}}, and it is scheduled to be presented at the next \emph{European Political Studies Association} (Prague) and the next \emph{Latin American Studies Association} (San Francisco) conferences. Taking these findings as a starting point, \href{https://www.utu.fi/en/people/janne-tukiainen}{Prof. Janne Tukiainen} (Turku School of Economics) and myself are currently collaborating in an improved experimental design that considers income inequality and political risk. We are interested in the estimated propensities of vote selling under different scenarios of poverty and income inequality---two rather different ways of ``aggregating'' economic deprivation, something that has not be studied in the experimental study of clientelism. For these matters, partnering with the \href{https://pcrclab.utu.fi/?page_id=894&lang=en}{PCRClab} will be fundamental to accomplish this goal. In addition to that, I was invited to collaborate with \href{https://www.utu.fi/en/people/outi-sarpila}{Prof. Outi Sarpila} (INVEST), where she and her research group are developing a series of experimental vignettes in the context of the \href{https://soma.utu.fi}{SOMA} project. Teaching-wise, I am also really excited to share that I've been scheduled to co-teach the ``\emph{Experimental Social Science}'' course next year for the INVEST doctoral and master students

% Research: Statistical
Quantitative/statistical work is another area that I also actively exploit in my research and in my teaching. For instance, after publishing the ``\emph{Inclusive Institutions, Unequal Outcomes}'' paper---where we used time-series methods---I started collaborating with \href{https://www.utu.fi/en/people/heikki-kauppi}{Prof. Heikki Kauppi} and \href{https://www.utu.fi/en/people/wilma-nissila}{Wilma Nissil{\"a}}, both from the Turku School of Economics. In this \href{https://github.com/hbahamonde/Bahamonde_Kovac/raw/master/abstract.txt}{paper} we try to introduce Global Vector Auto-regression models (GVAR) into political science. Since this is a joint effort with my coauthor from the University of Cincinnati, I believe this is also a nice opportunity to spread the word beyond Europe about our University and INVEST. Finally, I am currently collaborating in a very active way with \href{https://www.utu.fi/en/people/outi-sarpila}{Prof. Outi Sarpila} (INVEST) in a paper entitled ``\href{https://raw.githubusercontent.com/hbahamonde/Physical/main/abstract.rtf}{Physical Appearance and Turnout in Finnish Elections: An Inequality Perspective}.'' Here we argue that there exists a systematic electoral penalty, particularly for female candidates that look-like and also hold working-class occupations. We are currently applying to present this paper at the next \href{https://nsa2022.is}{Nordic Sociological Association} meeting in Reykjavik, Iceland. Teaching-wise, I'm beyond thrilled to also share that during March 2022 I am teaching the \emph{Quantitative Methods} module for the Masters in Social Sciences offered by the Faculty of Social Sciences. 

% in Experimental Research in Social and Behavioral Sciences
It is my belief that by holding the tenure-track position I will be able to boost up my current contributions to our University, Faculty of Social Science and INVEST. For example, I was recently appointed \emph{Research Director} of the \emph{INVEST-Hub} specialized in experimental research. This new position will allow me to organize different activities like seminars (or ``brown bags'') oriented to connect in-house researchers. By opening these activities to graduate students and faculty members of the University of Turku and \emph{beyond}, my idea is to position our researchers as leading scholars in the context of experimental social sciences. However, by holding the tenure-track position, I will be able to do more than that. For instance, supervising master theses and doctoral dissertations, specially in the context of their experimental designs. But also, mentoring undergraduate and young scholars in their earlier careers is also part of the service I would love to give to our community. 

% service
In terms of service, leadership and reach out, I also think that by holding the tenure-track position I will be able to contribute even more. For instance, my colleague \href{https://research.abo.fi/en/persons/inga-saikkonen}{Prof. Inga Saikkonen} (\AA bo Akademi) and myself, have put two panels together for the next \emph{European Consortium for Political Research} (Innsbruck, Austria) and the next \emph{Finnish Political Science Association} (Helsinki) meetings. I look forward to chairing these two panels. Also, I am an Editorial Board member of the \href{https://www.springer.com/journal/43545}{Social Sciences} journal (Springer). However, I strongly believe that holding the tenure-track position will allow me to go beyond these initiatives. For example, by being a professor, I might have more leverage to gather in-house researchers and graduate students to propose entire \emph{international} workshops and conferences on experimental social sciences, both on-campus \emph{and} overseas. The whole point will be to position our University and our scholars as a leading group in experimental research. I think that sending these kinds of credible signals to the general academic community, communicates the idea that the Faculty of Social Sciences and INVEST are vibrant and \emph{internationally}-oriented places to pursue research on experimental social sciences. It follows then that this should attract the attention of even more researchers and graduate students interested in pursuing these avenues of research. I'd be honored to contribute to the internationalization of our University.

% short plan: Academy of Finland
Should this opportunity be given to me, my plan for the next three to four years is to make sure I receive and implement an \emph{Academy of Finland} project. From an institutional-building point of view, I'd like to attract students and strengthen the usage of experiments in social sciences in our Faculty, the University and INVEST. From a substantive point of view, the project I am currently developing is about the relationship between income inequality and the electoral likability of populist candidates. In political science we are currently concerned about the widespread emergence of anti-system and populist candidates that get elected via traditional/electoral means. The danger this problem entails is that under these circumstances, democratic backslidings are very likely. Importantly, \emph{higher levels of income inequality offer one of the strongest explanations for the public's support for these kinds of candidates}. Methodologically speaking, my project seeks to develop a series of survey experiments to measure the public's support for populist candidates. Hopefully the project contributes to the problem by providing predicting tools for these kinds of behaviors. At the elite levels, the project intends to identify \emph{and} predict ``populist'' discourses by systematizing aggregated data (speeches and Twitter data) via artificial intelligence methods. In terms of the scope, my intention is to field this project in Finland, but then extend it to other countries in Europe as well as a sample of Latin American countries. Unfortunately, this phenomena is quickly spreading in all these mentioned cases. In fact, this project promises to be of high-impact because Finland might not be an exception to this phenomena. The impact goes beyond academic journals and university seminar rooms. In fact, I also expect to directly contribute to Finnish society not only by publicly speaking solo or with other experts about this problem, but also by appearing in different media outlets about this urgent and important topic.\footnote{I invite the interested reader to check out in my CV the Op-eds and media appearances I have had in the past while I was an assistant professor (tenure-track) in Chile. I must say that most of them are in Spanish, though.} Since I had a chance to \href{https://paloresearch.fi/en/events/}{discuss a paper} during the 2021 \href{https://paloresearch.fi/en/}{Participation in Long-Term Decision-Making} Conference, I already know that collaborating with \href{https://www.utu.fi/en/people/maija-setala}{Prof. Maija Set{\"a}l{\"a}} (Department of Philosophy, Contemporary History and Political Science, University of Turku) and the PALO project will be of utter importance. 

% to conclude
To conclude, it bears mentioning that I am beyond interested in this position because {\bf my main long-run professional priority is to build an academic career in the University of Turku.  Hence, my strategic goal is to get through the different evaluation processes that in time might lead to the tenured position. And for that, the first step in the ``ladder of tenure'' is to get into the tenure-track path}. In fact, I've been offered the possibility to start soon with the application process for the title of \emph{Docent}. For me, any element that links me up with the University of Turku is extremely valuable. 

% personal
Finally, Finland---and the City of Turku in particular---has been a welcoming, friendly, and overall a really nice place to call \emph{home} and to rise our two small children (5 and 3 years-old). We all have fell in love with its people, culture, language, weather and its connection with nature. In fact, not only my wife and me are taking a \emph{Finnish for Beginners} language course, but also Tobias and Olimpia are attending an all-Finnish language daycare since October 2021. Since my wife has German passport, the immigration process for us might facilitate things, specially now that she's looking for a job in Turku. All in all, we've taken the INVEST Senior Researcher  opportunity very seriously from the very beginning---professionally and family-wise. All these personal factors, for sure, constitute another really motivation to apply for the tenure-track position. 

Thank you for considering my application. I look forward to hearing from you. Should you have any questions about my application, please do contact me. I'll be happy to clarify any doubts you might have.




\vspace{7cm}

{\hspace{12cm}Hector Bahamonde}


%----------------------------------------------------------------------------------------

\end{newlfm}
\end{document}
